% TOP_PROBLEM: {"problemId": "problem.p-versus-pspace", "canonicalTitle": "P versus PSPACE", "releaseRank": 14, "primaryDomain": "theoretical_computer_science", "primaryDomainLabel": "Theoretical computer science", "displayStatus": "open", "releaseStatus": "open"}
% !TeX program = lualatex
% Definition and sources only; this file is not an LLM solution attempt.
\documentclass[11pt]{article}
\usepackage[margin=25mm]{geometry}
\usepackage{fontspec}
\setmainfont{Latin Modern Roman}
\usepackage{amsmath,amssymb,mathtools}
\usepackage{unicode-math}
\setmathfont{Latin Modern Math}
\usepackage{xurl,hyperref}
\hypersetup{colorlinks=true,urlcolor=blue}
\setcounter{secnumdepth}{0}
\setlength{\parindent}{0pt}
\setlength{\parskip}{5pt}
\setlength{\emergencystretch}{3em}
\begin{document}
\section*{\texorpdfstring{P versus PSPACE}{P versus PSPACE}}\label{14}
\subsection{Definitions and mathematical statement}
For an input of length $n$, time counts computation steps and space counts work-tape cells. Define
\[
 \mathsf P=\bigcup_{k\ge1}\mathsf{DTIME}(n^k),\qquad
 \mathsf{PSPACE}=\bigcup_{k\ge1}\mathsf{DSPACE}(n^k).
\]
Machines must halt and decide their language correctly on every input. Determine whether these classes coincide. An algorithm with polynomial space and exponential time does not establish membership in $\mathsf P$.
\par\smallskip\textit{Formulation and concept references:} \href{https://art.torvergata.it/bitstream/2108/32769/60437/ComputabilityDecidabilityComplexity.pdf}{[S1]}.

\subsection{Short English statement}
Can every decision problem solvable with polynomially much memory also be solved in polynomially much time?
\subsection{Sources}
\begin{itemize}
\item[S1]Walter S. Brainerd and Lawrence H. Landweber et al., Computability, Decidability, Complexity, third edition.
\url{https://art.torvergata.it/bitstream/2108/32769/60437/ComputabilityDecidabilityComplexity.pdf}.
\textit{Formulation source.}
\item[S2]Simulating Time With Square-Root Space — Ryan Williams, February 24, 2025.
\url{https://eccc.weizmann.ac.il/report/2025/017/download}.
\textit{Further reference; not a proof endorsement.}
\item[S3]Some Recent Developments in Space Complexity — R. Ryan Williams, MFCS 2026.
\url{https://drops.dagstuhl.de/entities/document/10.4230/LIPIcs.MFCS.2026.3}.
\textit{Further reference; not a proof endorsement.}
\item[S4]Some conditions implying if P=NP then P=PSPACE — Ismael Rodriguez (February 10, 2026).
\url{https://arxiv.org/abs/2602.10073}.
\textit{Further reference; not a proof endorsement.}
\end{itemize}
\end{document}
